\begin{statementbox}
	\textbf{\textcolor{CStatement}{条件}}

	\begin{description}[style=nextline, leftmargin=2.8em, labelsep=0.5em, itemsep=0.35em]
		\item[\textbf{\textcolor{CStatement}{条件 1（椭圆标准式）：}}]
		      椭圆方程为
		      \[
			      \frac{x^2}{a^2}+\frac{y^2}{b^2}=1,
		      \]
		      其中 \(a>b>0\)。

		\item[\textbf{\textcolor{CStatement}{条件 2（直线一般式）：}}]
		      直线方程为
		      \[
			      Ax+By+C=0,
		      \]
		      其中 \(A,B\) 不全为零。

		\item[\textbf{\textcolor{CStatement}{条件 3（辅助量定义）：}}]
		      记
		      \[
			      R=\sqrt{a^2A^2+b^2B^2}.
		      \]
	\end{description}

	\textbf{\textcolor{CStatement}{结论}}

	\begin{description}[style=nextline, leftmargin=2.8em, labelsep=0.5em, itemsep=0.45em]
		\item[\textbf{\textcolor{CStatement}{结论一（最值公式）：}}]
		      \textbf{\textcolor{CStatement}{直观描述：}}
		      椭圆上的点到直线的最远距离为
		      \[
			      \frac{|C|+R}{\sqrt{A^2+B^2}},
		      \]
		      最近距离需根据直线与椭圆是否相交分情况讨论。

		      \[
			      d_{\min}
			      =
			      \max\left\{0,\frac{|C|-R}{\sqrt{A^2+B^2}}\right\},
			      \qquad
			      d_{\max}
			      =
			      \frac{|C|+R}{\sqrt{A^2+B^2}}.
		      \]

		\item[\textbf{\textcolor{CStatement}{推广结论（位置判定）：}}]
		      \textbf{\textcolor{CStatement}{直观描述：}}
		      通过比较 \(|C|\) 与 \(R\)，可判断直线与椭圆的位置关系。

		      \[
			      \left.
			      \begin{aligned}
				      |C|>R & \Longleftrightarrow \text{直线与椭圆相离（无公共点）}, \\
				      |C|=R & \Longleftrightarrow \text{直线与椭圆相切},       \\
				      |C|<R & \Longleftrightarrow \text{直线与椭圆相交于两个点}.
			      \end{aligned}
			      \right.
		      \]
	\end{description}

	\textbf{\textcolor{CStatement}{取等条件：}}

	令
	\[
		P_+=\left(\frac{a^2A}{R},\frac{b^2B}{R}\right),
		\qquad
		P_-=-P_+.
	\]
	则 \(Ax+By\) 在椭圆上分别于 \(P_+\)、\(P_-\) 处取得最大值 \(R\) 与最小值 \(-R\)。

	\begin{itemize}[leftmargin=*, itemsep=0.35em, topsep=0.4em]
		\item 当 \(|C|>R\)（相离）时，最近距离和最远距离分别在两个法向支撑点 \(P_+\)、\(P_-\) 中取得，具体对应哪一个点取决于 \(C\) 的符号。

		\item 当 \(|C|\le R\)（相交或相切）时，\(d_{\min}=0\) 在直线与椭圆的交点（或切点）处取得；\(d_{\max}\) 仍在相应的法向支撑点处取得。

		\item 特别地，若 \(C>0\)，则 \(d_{\max}\) 在 \(P_+\) 处取得；若 \(C<0\)，则 \(d_{\max}\) 在 \(P_-\) 处取得；若 \(C=0\)，则 \(P_+\)、\(P_-\) 都可取得最大距离。
	\end{itemize}

\end{statementbox}